\MathBookSourcePage{92}
\section{Stokes 方程}\label{sec:stokes-equations}

考虑描述不可压缩黏性流体 $N$ 维运动的 Navier--Stokes 方程:
\begin{equation}\label{eq:stokes-navier-momentum-components}
\rho\left(
\frac{\partial u_i}{\partial t}
+\sum_{j=1}^{N}u_j\frac{\partial u_i}{\partial x_j}
\right)
-\sum_{j=1}^{N}\frac{\partial\sigma_{ij}}{\partial x_j}
=\rho f_i,
\qquad 1\le i\le N,
\end{equation}
\begin{equation}\label{eq:stokes-incompressibility-strain}
\operatorname{div}\mathbf{u}
=\sum_{i=1}^{N}D_{ii}(\mathbf{u})=0
\qquad\text{(incompressibility condition)},
\end{equation}
其中
\begin{equation}\label{eq:stokes-stress-tensor}
\sigma_{ij}=-P\delta_{ij}+2\mu D_{ij}(\mathbf{u}),
\qquad 1\le i,j\le N,
\end{equation}
并且
\[
D_{ij}(\mathbf{u})
=\frac12\left(
\frac{\partial u_i}{\partial x_j}
+\frac{\partial u_j}{\partial x_i}
\right).
\]

\MathBookSourcePage{93}
在这些方程中, $\mathbf{u}=(u_1,\ldots,u_N)$ 是流体速度,
$\rho$ 是其密度(假设为常数), $\mu>0$ 是其黏度(也假设为常数),
$P$ 是其压力; $(\sigma_{ij})$ 是应力张量,
$\mathbf{f}=(f_1,\ldots,f_N)$ 表示单位质量的体力密度(例如重力).

与通常一样, 置
\begin{equation}\label{eq:stokes-kinematic-pressure-viscosity}
p=P/\rho,\qquad \nu=\mu/\rho.
\end{equation}
这里 $p$ 是运动压力, $\nu$ 是运动黏度, 但为简单起见,
下文仍将它们称为压力和黏度. 采用这些记号后, Navier--Stokes 方程变为
\begin{equation}\label{eq:stokes-navier-components}
\left\{
\begin{aligned}
\frac{\partial u_i}{\partial t}
+\sum_{j=1}^{N}u_j\frac{\partial u_i}{\partial x_j}
-2\nu\sum_{j=1}^{N}\frac{\partial D_{ij}(\mathbf{u})}{\partial x_j}
+\frac{\partial p}{\partial x_i}
&=f_i,\qquad 1\le i\le N,\\
\operatorname{div}\mathbf{u}&=0.
\end{aligned}
\right.
\end{equation}
注意, 当 $\operatorname{div}\mathbf{u}=0$ 时, 有恒等式
\begin{equation}\label{eq:stokes-strain-laplacian-identity}
\sum_{j=1}^{N}\frac{\partial D_{ij}(\mathbf{u})}{\partial x_j}
=\frac12\sum_{j=1}^{N}\left(
\frac{\partial^2u_i}{\partial x_j^2}
+\frac{\partial^2u_j}{\partial x_i\partial x_j}
\right)
=\frac12\Delta u_i,
\end{equation}
所以 \eqref{eq:stokes-navier-components} 可更方便地写成
\begin{equation}\label{eq:stokes-navier-vector}
\left\{
\begin{aligned}
\frac{\partial\mathbf{u}}{\partial t}
+\sum_{j=1}^{N}u_j\frac{\partial\mathbf{u}}{\partial x_j}
-\nu\Delta\mathbf{u}+\operatorname{grad}p&=\mathbf{f},\\
\operatorname{div}\mathbf{u}&=0.
\end{aligned}
\right.
\end{equation}

引入参数 Reynolds 数, 用它度量黏度对流动的影响.
对于一个给定问题, 设 $L$ 是特征长度, $U$ 是特征速度.
这确定了特征时间 $T=L/U$. 然后引入无量纲量
\[
x'=x/L,\qquad \mathbf{u}'=\mathbf{u}/U,\qquad t'=t/T.
\]
采用这一变量变换, 容易检验 Navier--Stokes 方程变为(采用显然的记号)
\[
\frac{\partial\mathbf{u}'}{\partial t'}
+\sum_{j=1}^{N}u'_j\frac{\partial\mathbf{u}'}{\partial x'_j}
-\frac{\nu}{LU}\Delta'\mathbf{u}'
+\operatorname{grad}'p'=\mathbf{f}',
\qquad
\operatorname{div}'\mathbf{u}'=0,
\]
其中
\[
p'=P/(\rho U^2),\qquad \mathbf{f}'=L\mathbf{f}/U^2.
\]

\MathBookSourcePage{94}
现在把 Reynolds 数定义为无量纲数
\[
\operatorname{Re}=LU/\nu.
\]
于是 Navier--Stokes 方程可用无量纲变量写成
\begin{equation}\label{eq:stokes-navier-dimensionless}
\left\{
\begin{aligned}
\frac{\partial\mathbf{u}}{\partial t}
+\sum_{j=1}^{N}u_j\frac{\partial\mathbf{u}}{\partial x_j}
-\frac1{\operatorname{Re}}\Delta\mathbf{u}
+\operatorname{grad}p&=\mathbf{f},\\
\operatorname{div}\mathbf{u}&=0.
\end{aligned}
\right.
\end{equation}
这再次得到 \eqref{eq:stokes-navier-vector}, 只是把 $\nu$ 换成了
$1/\operatorname{Re}$.

暂且对 \eqref{eq:stokes-navier-components} 或
\eqref{eq:stokes-navier-vector} 作两点简化. 我们只考虑稳态(或定常)情形,
即 $\partial\mathbf{u}/\partial t=\mathbf{0}$; 此外, 假设速度
$\mathbf{u}$ 足够小, 从而可以忽略非线性对流项
$u_j(\partial u_i/\partial x_j)$. 于是得到 Stokes 方程
\begin{equation}\label{eq:stokes-stationary-strain}
\left\{
\begin{aligned}
-2\nu\sum_{j=1}^{N}
\frac{\partial D_{ij}(\mathbf{u})}{\partial x_j}
+\frac{\partial p}{\partial x_i}
&=f_i,\qquad 1\le i\le N,\\
\operatorname{div}\mathbf{u}&=0,
\end{aligned}
\right.
\end{equation}
它可更方便地写成
\begin{equation}\label{eq:stokes-stationary-vector}
\left\{
\begin{aligned}
-\nu\Delta\mathbf{u}+\operatorname{grad}p&=\mathbf{f},\\
\operatorname{div}\mathbf{u}&=0.
\end{aligned}
\right.
\end{equation}
Stokes 方程是线性的, 但由于不可压缩条件
$\operatorname{div}\mathbf{u}=0$, 它们仍值得特别关注.
本节将证明 Stokes 方程解的存在唯一性, 并导出若干变分形式,
供以后作逼近时使用.

\subsection{速度--压力形式的 Dirichlet 问题}
\label{subsec:stokes-dirichlet-velocity-pressure}

为了使 Stokes 方程 \eqref{eq:stokes-stationary-vector} 构成适定问题,
必须为其补充适当的边界条件. 先考虑 Dirichlet 边界条件.

\setcounter{statement}{0}
\begin{Theorem}\label{thm:stokes-dirichlet-well-posedness}
设 $\Omega$ 是 $\mathbb{R}^N$ 中有界连通开集, 其边界
$\Gamma$ 是 Lipschitz 连续的. 给定
$\mathbf{f}\in H^{-1}(\Omega)^N$ 和
$\mathbf{g}\in H^{1/2}(\Gamma)^N$, 且
\begin{equation}\label{eq:stokes-boundary-flux-compatibility}
\int_\Gamma\mathbf{g}\cdot\mathbf{n}\,\mathrm{d}s=0,
\end{equation}
则存在唯一一对
$(\mathbf{u},p)\in H^1(\Omega)^N\times L_0^2(\Omega)$,
它是方程
\begin{equation}\label{eq:stokes-dirichlet-problem}
\left\{
\begin{aligned}
-\nu\Delta\mathbf{u}+\operatorname{grad}p&=\mathbf{f}
&&\text{in }\Omega,\\
\operatorname{div}\mathbf{u}&=0
&&\text{in }\Omega,\\
\mathbf{u}&=\mathbf{g}
&&\text{on }\Gamma
\end{aligned}
\right.
\end{equation}
的解.
\end{Theorem}
\begin{Proof}
由 \eqref{eq:stokes-boundary-flux-compatibility} 和 Lemma
\ref{lem:sec2-div-free-lifting}, 存在函数
$\mathbf{u}_0\in H^1(\Omega)^N$, 使
\[
\operatorname{div}\mathbf{u}_0=0\quad\text{in }\Omega,
\qquad
\mathbf{u}_0=\mathbf{g}\quad\text{on }\Gamma.
\]

\MathBookSourcePage{95}
现在把问题 \eqref{eq:stokes-dirichlet-problem} 纳入 Section
\ref{sec:abstract-variational-problem} 的框架. 置
\begin{equation}\label{eq:stokes-dirichlet-abstract-data}
\begin{gathered}
X=H_0^1(\Omega)^N,\qquad M=L_0^2(\Omega),\\
\|\cdot\|_X=|\cdot|_{1,\Omega},
\qquad
\|\cdot\|_M=\|\cdot\|_{0,\Omega},\\
a(\mathbf{u},\mathbf{v})
=\nu\sum_{i,j=1}^{N}
\left(
\frac{\partial u_i}{\partial x_j},
\frac{\partial v_i}{\partial x_j}
\right)
=\nu(\operatorname{grad}\mathbf{u},
\operatorname{grad}\mathbf{v}),\\
b(\mathbf{v},q)=-(q,\operatorname{div}\mathbf{v}),\\
\langle\mathbf{l},\mathbf{v}\rangle
=\langle\mathbf{f},\mathbf{v}\rangle
-a(\mathbf{u}_0,\mathbf{v}),
\qquad \chi=0.
\end{gathered}
\end{equation}
于是
\[
V=\{\mathbf{v}\in H_0^1(\Omega)^N;
\ \operatorname{div}\mathbf{v}=0\}.
\]
必须检验形式 $a(\cdot,\cdot)$ 是 $V$-椭圆的,
并且形式 $b(\cdot,\cdot)$ 满足 inf-sup 条件
\eqref{eq:abstract-inf-sup}. 一方面, 椭圆性显然, 因为
\[
a(\mathbf{v},\mathbf{v})=\nu|\mathbf{v}|_{1,\Omega}^2.
\]
另一方面, inf-sup 条件是
\begin{equation}\label{eq:stokes-pressure-infsup}
\sup_{\mathbf{v}\in H_0^1(\Omega)^N}
\frac{(q,\operatorname{div}\mathbf{v})}
{|\mathbf{v}|_{1,\Omega}}
\ge\beta\|q\|_{0,\Omega}
\qquad\forall q\in L_0^2(\Omega).
\end{equation}
设 $q\in L_0^2(\Omega)$. 由 Corollary
\ref{cor:sec2-grad-div-isomorphisms}, 存在唯一函数
$\mathbf{v}\in V^\perp$, 使
\[
\operatorname{div}\mathbf{v}=q,
\qquad
|\mathbf{v}|_{1,\Omega}\le C\|q\|_{0,\Omega}.
\]
因此
\[
\frac{(q,\operatorname{div}\mathbf{v})}
{|\mathbf{v}|_{1,\Omega}}
=\frac{\|q\|_{0,\Omega}^2}
{|\mathbf{v}|_{1,\Omega}}
\ge\frac1C\|q\|_{0,\Omega},
\]
由此以 $\beta=1/C$ 得到 \eqref{eq:stokes-pressure-infsup}.

现在可以应用 Corollary
\ref{cor:abstract-v-elliptic-well-posedness}: 存在唯一一对函数
$(\mathbf{w},p)\in H_0^1(\Omega)^N\times L_0^2(\Omega)$, 使
\[
\left\{
\begin{aligned}
a(\mathbf{w},\mathbf{v})+b(\mathbf{v},p)
&=\langle\mathbf{l},\mathbf{v}\rangle
&&\forall\mathbf{v}\in H_0^1(\Omega)^N,\\
b(\mathbf{w},q)&=0
&&\forall q\in L_0^2(\Omega).
\end{aligned}
\right.
\]

\MathBookSourcePage{96}
等价地,
$(\mathbf{u}=\mathbf{u}_0+\mathbf{w},p)
\in[\mathbf{u}_0+H_0^1(\Omega)^N]\times L_0^2(\Omega)$
是方程
\[
\left\{
\begin{aligned}
\nu(\operatorname{grad}\mathbf{u},
\operatorname{grad}\mathbf{v})
-(p,\operatorname{div}\mathbf{v})
&=\langle\mathbf{f},\mathbf{v}\rangle
&&\forall\mathbf{v}\in H_0^1(\Omega)^N,\\
(q,\operatorname{div}\mathbf{u})&=0
&&\forall q\in L_0^2(\Omega)
\end{aligned}
\right.
\]
的解. 再次应用 Corollary
\ref{cor:sec2-grad-div-isomorphisms}, 最后一个方程等价于
$\operatorname{div}\mathbf{u}=0$. 此外,
$\mathbf{u}\in\mathbf{u}_0+H_0^1(\Omega)^N$ 当且仅当
\[
\mathbf{u}\in H^1(\Omega)^N,
\qquad
\mathbf{u}|_\Gamma=\mathbf{g}.
\]
因此存在唯一一对
$(\mathbf{u},p)\in H^1(\Omega)^N\times L_0^2(\Omega)$, 使
\[
\left\{
\begin{aligned}
\nu(\operatorname{grad}\mathbf{u},
\operatorname{grad}\mathbf{v})
-(p,\operatorname{div}\mathbf{v})
&=\langle\mathbf{f},\mathbf{v}\rangle
&&\forall\mathbf{v}\in H_0^1(\Omega)^N,\\
\operatorname{div}\mathbf{u}&=0,\\
\mathbf{u}|_\Gamma&=\mathbf{g}.
\end{aligned}
\right.
\]
现在利用经典论证, 容易证明这一问题等价于问题
\eqref{eq:stokes-dirichlet-problem}.
\end{Proof}

\setcounter{statement}{0}
\begin{Remark}\label{rem:stokes-pressure-quotient-space}
\eqref{eq:stokes-dirichlet-abstract-data} 中选择
$M=L_0^2(\Omega)$ 只是为了方便, 也完全可以取
$M=L^2(\Omega)/\mathbb{R}$. 另一方面, 在上述证明中也可以选择
\[
a(\mathbf{u},\mathbf{v})
=2\nu\sum_{i,j=1}^{N}
\int_\Omega D_{ij}(\mathbf{u})
D_{ij}(\mathbf{v})\,\mathrm{d}x.
\]
这是下列恒等式的结果:
\begin{equation}\label{eq:stokes-strain-gradient-identity}
\sum_{i,j=1}^{N}
\int_\Omega D_{ij}(\mathbf{u})
D_{ij}(\mathbf{v})\,\mathrm{d}x
=\frac12\sum_{i,j=1}^{N}
\int_\Omega
\frac{\partial u_i}{\partial x_j}
\frac{\partial v_i}{\partial x_j}\,\mathrm{d}x,
\end{equation}
该恒等式对所有满足
$\operatorname{div}\mathbf{u}=0$ 的
$\mathbf{u}\in H^1(\Omega)^N$ 以及所有
$\mathbf{v}\in H_0^1(\Omega)^N$ 成立. 事实上, 利用算子
$D_{ij}$ 关于 $i,j$ 的对称性, 有
\[
\sum_{i,j=1}^{N}\int_\Omega
D_{ij}(\mathbf{u})D_{ij}(\mathbf{v})\,\mathrm{d}x
=\sum_{i,j=1}^{N}\int_\Omega
D_{ij}(\mathbf{u})
\frac{\partial v_i}{\partial x_j}\,\mathrm{d}x.
\]
此外,
\[
\sum_{i,j=1}^{N}\int_\Omega
\frac{\partial u_j}{\partial x_i}
\frac{\partial v_i}{\partial x_j}\,\mathrm{d}x
=-\sum_{i=1}^{N}
\left\langle
\frac{\partial\operatorname{div}\mathbf{u}}{\partial x_i},
v_i
\right\rangle
=0,
\]
由此得到 \eqref{eq:stokes-strain-gradient-identity}.
\end{Remark}
